% ============================================================
% ANÁLISE DO MAPA TAXONÔMICO — 6 EUF (2011-2016) + 3 ENA (2024-2026)
% Raciocínios Científicos (Taxonomia 350) + Bloom + Correlações
% Dataset: 150 questões (60 EUF + 90 ENA) | Mestrado PPGTE/CT/UFC
% ============================================================
\documentclass[12pt,a4paper]{article}
\usepackage[utf8]{inputenc}
\usepackage[T1]{fontenc}
\usepackage{lmodern}
\usepackage[brazilian]{babel}
\usepackage{geometry}
\geometry{margin=2.5cm,top=2.5cm,bottom=2cm}
\usepackage{indentfirst}
\usepackage{setspace}
\onehalfspacing
\usepackage{amsmath,amssymb}
\usepackage{booktabs,array,longtable}
\usepackage{hyperref}
\usepackage{graphicx}
\usepackage{float}
\usepackage{xcolor}
\usepackage{caption}

\hypersetup{colorlinks=true,linkcolor=blue,citecolor=blue,urlcolor=blue}

\definecolor{azulprof}{RGB}{0,51,102}
\definecolor{verdeexame}{RGB}{0,100,50}
\definecolor{vermelho}{RGB}{180,30,30}

\newcommand{\rcode}[1]{\textcolor{azulprof}{\textbf{#1}}}
\newcommand{\bloom}[1]{\textcolor{verdeexame}{\textbf{#1}}}
\newcommand{\cat}[1]{\textcolor{vermelho}{\textbf{#1}}}

\title{\textbf{Análise do Mapa Taxonômico:}\\
\large 150 Questões (EUF 2011--2016 + ENA 2024--2026) \\
segundo a Taxonomia 350 de Raciocínios Científicos \\
e a Taxonomia de Bloom}
\author{OpenCode Ecosystem --- Análise de Dados\\
Mestrado PPGTE/CT/UFC}
\date{Maio 2026}

\begin{document}
\maketitle

\begin{abstract}
Este relatório apresenta a análise completa do mapeamento taxonômico de
150 questões: 60 dos 6 Exames Unificados de Física (EUF, 2011--2016) e
90 dos 3 Exames Nacionais de Matemática (ENA, 2024--2026).
Cada questão foi classificada segundo: (i) a Taxonomia 350 de Raciocínios
Científicos (R-codes primário e secundário), (ii) a categoria taxonômica
correspondente (5 categorias da Taxonomia 350), e (iii) o nível de Bloom
(Aplicar ou Analisar). São analisadas as correlações entre área,
R-code e nível de Bloom, com testes de $\chi^2$ e V de Cramer, tanto
separadamente por exame quanto de forma agregada.
\vspace{0.3cm}

\noindent\textbf{Palavras-chave}: EUF. ENA. Taxonomia 350. Raciocínio científico.
Classificação de questões. Bloom. Correlações.
\end{abstract}

\tableofcontents
\newpage

% ============================================================
\section{Visão Geral do Dataset}
% ============================================================

O dataset completo é composto por 150 questões, das quais 60 são do
tipo EUF (6 exames $\times$ 10 questões) e 90 do tipo ENA
(3 exames $\times$ 30 questões). Cada questão da EUF cobre
uma das 6 áreas da física: Eletromagnetismo, Física Moderna, Mecânica
Clássica, Mecânica Estatística, Mecânica Quântica e Termodinâmica.
As 90 questões do ENA cobrem 16 áreas da matemática, incluindo
Álgebra, Geometria, Teoria dos Números, Estatística, entre outras.

\subsection{Distribuição por Exame}

A Tabela~\ref{tab:exames} apresenta os 6 exames analisados.

\begin{table}[H]
\centering
\caption{Exames EUF no dataset}
\label{tab:exames}
\begin{tabular}{lccc}
\toprule
\textbf{Exame} & \textbf{Questões} & \textbf{R-codes Primários} & \textbf{R-codes Total} \\
\midrule
EUF 2011-1 & 10 & 7 & 8 \\
EUF 2012-2 & 10 & 7 & 8 \\
EUF 2014-2 & 10 & 7 & 9 \\
EUF 2015-1 & 10 & 5 & 6 \\
EUF 2015-2 & 10 & 7 & 8 \\
EUF 2016-1 & 10 & 8 & 9 \\
\midrule
\textbf{Total} & \textbf{60} & \textbf{15} & \textbf{20} \\
\bottomrule
\end{tabular}
\end{table}

\subsection{Distribuição por Área da Física}

Cada área da física aparece 10 vezes no dataset (6 exames, sendo 2 questões
de Eletromagnetismo, 2 de Física Moderna, 2 de Mecânica Clássica, 1 de
Mecânica Estatística, 2 de Mecânica Quântica e 1 de Termodinâmica por
exame), totalizando 12 questões para as áreas com 2 questões por exame
e 6 para as áreas com 1 questão por exame.

% ============================================================
\section{Classificação Taxonômica dos R-codes}
% ============================================================

Os 20 R-codes distintos identificados no dataset foram mapeados para
5 categorias da Taxonomia 350. A Tabela~\ref{tab:mapa_completo}
apresenta o mapeamento completo.

\begin{longtable}{@{}p{2.2cm}p{3.8cm}p{3cm}p{4.5cm}@{}}
\caption{Mapa completo: 20 R-codes $\rightarrow$ Categorias Taxonomia 350}
\label{tab:mapa_completo}\\
\toprule
\textbf{Código} & \textbf{Descrição} & \textbf{Categoria} & \textbf{Ocorrências (P/S)} \\
\midrule
\endfirsthead

\caption[]{Mapa completo (continuação)}\\
\toprule
\textbf{Código} & \textbf{Descrição} & \textbf{Categoria} & \textbf{Ocorrências (P/S)} \\
\midrule
\endhead

\bottomrule
\endfoot

% CAT I
\rcode{R008} & Deducao-Formal-Passo-a-Passo
  & I. Lógico-Dedutivos & 1P + 1S \\
\rcode{R011} & Inducao-Estrutural
  & I. Lógico-Dedutivos & 1P + 0S \\
\midrule

% CAT II
\rcode{R010} & Decomposicao-Modular
  & II. Analítico-Estruturais & 1P + 0S \\
\midrule

% CAT VII
\rcode{R054} & Otimizacao-Restrita
  & VII. Econômico-Decisórios & 0P + 1S \\
\midrule

% CAT XIII
\rcode{R035} & Analise-Grafica-Experimental
  & XIII. Físico-Matemáticos & 0P + 0S (ENA) \\
\rcode{R065} & Simetria-Conservacao
  & XIII. Físico-Matemáticos & 8P + 3S \\
\rcode{R066} & Dimensional-Analitico
  & XIII. Físico-Matemáticos & 0P + 1S \\
\rcode{R067} & Perturbativo
  & XIII. Físico-Matemáticos & 1P + 0S \\
\rcode{R069} & Variacional-Principio
  & XIII. Físico-Matemáticos & 8P + 6S \\
\rcode{R070} & Limite-Assintotico
  & XIII. Físico-Matemáticos & 1P + 2S \\
\rcode{R071} & Dualidade-Transformacao
  & XIII. Físico-Matemáticos & 10P + 2S \\
\rcode{R072} & Estabilidade-Lyapunov
  & XIII. Físico-Matemáticos & 0P + 0S (ENA) \\
\rcode{R073} & Caos-Deterministico
  & XIII. Físico-Matemáticos & 1P + 0S \\
\rcode{R114} & Probabilidade-Bayesiana/Contagem
  & XIII. Físico-Matemáticos & 0P + 0S (ENA) \\
\rcode{R117} & Monte-Carlo-MCMC
  & XIII. Físico-Matemáticos & 9P + 3S \\
\rcode{R183} & Mecanica-Quantica-Particulas
  & XIII. Físico-Matemáticos & 11P + 0S \\
\rcode{R201} & Superposicao-Quantica
  & XIII. Físico-Matemáticos & 1P + 0S \\
\rcode{R202} & Emaranhamento-Quantico
  & XIII. Físico-Matemáticos & 2P + 0S \\
\midrule

% CAT XIV
\rcode{R109} & Espectral-Operadores
  & XIV. Lógica Matemática Avançada & 1P + 2S \\
\rcode{R229} & Homotopia-Fundamental
  & XIV. Lógica Matemática Avançada & 4P + 1S \\
\bottomrule
\end{longtable}

\noindent Onde: P = ocorrência como raciocínio primário, S = ocorrência
como raciocínio secundário, ENA = presente apenas em questões ENA
(não-EUF), não contabilizado nas 60 questões EUF.

\subsection{Distribuição Agregada por Categoria}

A Tabela~\ref{tab:categorias_agregadas} apresenta a distribuição das
ocorrências de raciocínio primário pelas 5 categorias.

\begin{table}[H]
\centering
\caption{Distribuição das ocorrências primárias por categoria}
\label{tab:categorias_agregadas}
\begin{tabular}{lcc}
\toprule
\textbf{Categoria} & \textbf{Ocorrências (P)} & \textbf{Percentual} \\
\midrule
XIII. Físico-Matemáticos & 52 & 86,7\% \\
XIV. Lógica Matemática Avançada & 5 & 8,3\% \\
I. Lógico-Dedutivos & 2 & 3,3\% \\
II. Analítico-Estruturais & 1 & 1,7\% \\
VII. Econômico-Decisórios & 0 & 0,0\% \\
\midrule
\textbf{Total} & \textbf{60} & \textbf{100\%} \\
\bottomrule
\end{tabular}
\end{table}

A Categoria XIII (Físico-Matemáticos) domina esmagadoramente com 86,7\%
das ocorrências primárias. Este resultado era esperado, dado que o EUF
é um exame de física teórica que demanda intensamente raciocínio
quantitativo-matemático.

\subsection{R-codes por Categoria nos 6 Exames}

A Tabela~\ref{tab:cat_por_exame} mostra a predominância da Cat.~XIII
em todos os exames, com a EUF 2014-2 apresentando a maior diversidade
categórica (4 categorias).

\begin{table}[H]
\centering
\caption{Categorias por exame (ocorrências primárias)}
\label{tab:cat_por_exame}
\begin{tabular}{lcccccc}
\toprule
\textbf{Categoria} & \textbf{2011-1} & \textbf{2012-2} & \textbf{2014-2} & \textbf{2015-1} & \textbf{2015-2} & \textbf{2016-1} \\
\midrule
I. Lógico-Dedutivos & 0 & 0 & 1 & 0 & 1 & 0 \\
II. Analítico-Estruturais & 0 & 0 & 1 & 0 & 0 & 0 \\
VII. Econômico-Decisórios & 0 & 0 & 0 & 0 & 0 & 0 \\
XIII. Físico-Matemáticos & 9 & 9 & 7 & 9 & 9 & 9 \\
XIV. Lógica Matemática Avançada & 1 & 1 & 1 & 1 & 0 & 1 \\
\midrule
\textbf{Total} & \textbf{10} & \textbf{10} & \textbf{10} & \textbf{10} & \textbf{10} & \textbf{10} \\
\bottomrule
\end{tabular}
\end{table}

% ============================================================
\section{Distribuição dos Níveis de Bloom}
% ============================================================

As 60 questões EUF foram classificadas em dois níveis de Bloom:
\bloom{Aplicar} (37 ocorrências) e \bloom{Analisar} (23 ocorrências).

\subsection{Bloom por Exame}

A Tabela~\ref{tab:bloom_exame} mostra a distribuição consistente de
Bloom em todos os exames.

\begin{table}[H]
\centering
\caption{Bloom por exame EUF}
\label{tab:bloom_exame}
\begin{tabular}{lccc}
\toprule
\textbf{Exame} & \textbf{Aplicar} & \textbf{Analisar} & \textbf{Total} \\
\midrule
EUF 2011-1 & 6 & 4 & 10 \\
EUF 2012-2 & 6 & 4 & 10 \\
EUF 2014-2 & 6 & 4 & 10 \\
EUF 2015-1 & 6 & 4 & 10 \\
EUF 2015-2 & 7 & 3 & 10 \\
EUF 2016-1 & 6 & 4 & 10 \\
\midrule
\textbf{Total} & \textbf{37} & \textbf{23} & \textbf{60} \\
\bottomrule
\end{tabular}
\end{table}

A predominância de \bloom{Aplicar} (61,7\%) indica que a EUF privilegia
a aplicação de métodos consagrados a problemas específicos. Apenas a
EUF 2015-2 foge do padrão 6A/4AN com 7A/3AN.

\subsection{Bloom por Área da Física}

A Tabela~\ref{tab:bloom_area} cruza Bloom com área, revelando que
Eletromagnetismo e Mecânica Clássica concentram mais \bloom{Aplicar},
enquanto Termodinâmica e Física Moderna têm maior proporção de
\bloom{Analisar}.

\begin{table}[H]
\centering
\caption{Bloom por área (todas as 60 EUF)}
\label{tab:bloom_area}
\begin{tabular}{lccc}
\toprule
\textbf{Área} & \textbf{Aplicar} & \textbf{Analisar} & \textbf{Total} \\
\midrule
Eletromagnetismo & 10 & 2 & 12 \\
Física Moderna & 4 & 8 & 12 \\
Mecânica Clássica & 10 & 2 & 12 \\
Mecânica Estatística & 4 & 2 & 6 \\
Mecânica Quântica & 6 & 6 & 12 \\
Termodinâmica & 3 & 3 & 6 \\
\midrule
\textbf{Total} & \textbf{37} & \textbf{23} & \textbf{60} \\
\bottomrule
\end{tabular}
\end{table}

% ============================================================
\section{Inclusão dos Exames ENA (Matemática)}
% ============================================================

O banco de dados foi expandido com 90 questões dos Exames Nacionais de
Matemática (ENA) dos anos 2024, 2025 e 2026, cobrindo 16 áreas da
matemática (Álgebra, Álgebra Linear, Aritmética, Combinatória,
Estatística, Geometria, Geometria Espacial, Geometria Plana, Lógica,
Matemática Financeira, Otimização, Probabilidade, Sequências,
Teoria dos Conjuntos, Teoria dos Números e Trigonometria).

\subsection{Distribuição de Bloom nos Exames ENA}

A Tabela~\ref{tab:bloom_ena_exame} mostra a distribuição consistente de
Bloom nos três anos, com pequena predominância de \bloom{Aplicar} (55,6\%)
sobre \bloom{Analisar} (44,4\%).

\begin{table}[H]
\centering
\caption{Bloom por exame ENA}
\label{tab:bloom_ena_exame}
\begin{tabular}{lccc}
\toprule
\textbf{Exame} & \textbf{Aplicar} & \textbf{Analisar} & \textbf{Total} \\
\midrule
ENA 2024 & 16 & 14 & 30 \\
ENA 2025 & 17 & 13 & 30 \\
ENA 2026 & 17 & 13 & 30 \\
\midrule
\textbf{Total} & \textbf{50} & \textbf{40} & \textbf{90} \\
\bottomrule
\end{tabular}
\end{table}

\subsection{Bloom ENA por Área da Matemática}

A Tabela~\ref{tab:bloom_ena_area} revela contrastes entre as áreas:
Aritmética e Geometria Plana concentram a maior parte das questões de
\bloom{Aplicar}, enquanto Álgebra, Teoria dos Números e Lógica são
predominantemente \bloom{Analisar}.

\begin{table}[H]
\centering
\caption{Bloom ENA por área da matemática}
\label{tab:bloom_ena_area}
\begin{tabular}{lccc}
\toprule
\textbf{Área} & \textbf{Aplicar} & \textbf{Analisar} & \textbf{Total} \\
\midrule
Álgebra & 6 & 11 & 17 \\
Álgebra Linear & 2 & 4 & 6 \\
Aritmética & 8 & 0 & 8 \\
Combinatória & 5 & 3 & 8 \\
Estatística & 0 & 3 & 3 \\
Geometria & 1 & 0 & 1 \\
Geometria Espacial & 3 & 0 & 3 \\
Geometria Plana & 15 & 1 & 16 \\
Lógica & 0 & 3 & 3 \\
Matemática Financeira & 1 & 0 & 1 \\
Otimização & 0 & 1 & 1 \\
Probabilidade & 4 & 0 & 4 \\
Sequências & 4 & 3 & 7 \\
Teoria dos Conjuntos & 0 & 1 & 1 \\
Teoria dos Números & 1 & 7 & 8 \\
Trigonometria & 0 & 3 & 3 \\
\midrule
\textbf{Total} & \textbf{50} & \textbf{40} & \textbf{90} \\
\bottomrule
\end{tabular}
\end{table}

\subsection{Bloom ENA por R-code}

A Tabela~\ref{tab:bloom_ena_rcode} detalha a relação entre os R-codes
identificados no ENA e o nível de Bloom.

\begin{table}[H]
\centering
\caption{Bloom ENA por R-code primário}
\label{tab:bloom_ena_rcode}
\begin{tabular}{lccc}
\toprule
\textbf{R-code} & \textbf{Aplicar} & \textbf{Analisar} & \textbf{Total} \\
\midrule
R008 (Deducao-Formal) & 32 & 13 & 45 \\
R009 (Critica-por-Contraexemplo) & 0 & 4 & 4 \\
R012 (Generalizacao-Estatistica) & 2 & 3 & 5 \\
R014 (Generalizacao-Matematica) & 0 & 8 & 8 \\
R020 (Construcao-Escalonada) & 5 & 3 & 8 \\
R023 (Reducao-ao-Absurdo) & 0 & 4 & 4 \\
R048 (Custo-Beneficio) & 1 & 0 & 1 \\
R054 (Otimizacao-Restrita) & 0 & 3 & 3 \\
R065 (Simetria-Conservacao) & 5 & 1 & 6 \\
R066 (Dimensional-Analitico) & 3 & 0 & 3 \\
R101 (Topologico-Ponto-Fixo) & 0 & 1 & 1 \\
R114 (Probabilidade-Contagem) & 2 & 0 & 2 \\
\midrule
\textbf{Total} & \textbf{50} & \textbf{40} & \textbf{90} \\
\bottomrule
\end{tabular}
\end{table}

Observa-se que R-codes como \rcode{R009} (Critica-por-Contraexemplo),
\rcode{R014} (Generalizacao-Matematica) e \rcode{R023}
(Reducao-ao-Absurdo) aparecem exclusivamente como \bloom{Analisar},
exigindo raciocínio crítico e abstrato. Já \rcode{R008}
(Deducao-Formal), o mais frequente (45 ocorrências, 50\% do ENA),
distribui-se majoritariamente como \bloom{Aplicar} (71\%).

\subsection{Comparação EUF versus ENA}

A Tabela~\ref{tab:euf_ena_comparativo} sintetiza as diferenças entre
os dois tipos de exame.

\begin{table}[H]
\centering
\caption{Comparação EUF versus ENA}
\label{tab:euf_ena_comparativo}
\begin{tabular}{lcc}
\toprule
\textbf{Dimensão} & \textbf{EUF (Física)} & \textbf{ENA (Matemática)} \\
\midrule
Questões & 60 & 90 \\
R-codes distintos & 15 & 12 \\
R-codes em comum & R008, R065 & R008, R065 \\
Categorias Tax.~350 & 5 & 5 \\
Bloom Aplicar & 37 (61,7\%) & 50 (55,6\%) \\
Bloom Analisar & 23 (38,3\%) & 40 (44,4\%) \\
Cat.~dominante & XIII (86,7\%) & I (67,8\%) \\
R-code mais frequente & R183 (11) & R008 (45) \\
\bottomrule
\end{tabular}
\end{table}

A diferença fundamental está na categoria dominante: enquanto o EUF
mobiliza majoritariamente raciocínios \textbf{Físico-Matemáticos}
(Cat.~XIII, 86,7\%), o ENA concentra-se em raciocínios
\textbf{Lógico-Dedutivos} (Cat.~I, 67,8\%). Esta divergência reflete
a natureza dos exames: a física teórica exige intenso raciocínio
quantitativo-matemático, enquanto a matemática pura privilegia a
dedução formal passo a passo.

\subsection{Distribuição Consolidada (150 Questões)}

Com a inclusão do ENA, a distribuição por categoria para o total de
150 questões é:

\begin{table}[H]
\centering
\caption{Categorias Taxonomia 350 --- Total de 150 questões}
\label{tab:categorias_150}
\begin{tabular}{lcc}
\toprule
\textbf{Categoria} & \textbf{Ocorrências} & \textbf{Percentual} \\
\midrule
I. Lógico-Dedutivos & 63 & 42,0\% \\
XIII. Físico-Matemáticos & 68 & 45,3\% \\
II. Analítico-Estruturais & 9 & 6,0\% \\
XIV. Lógica Matemática Avançada & 6 & 4,0\% \\
VII. Econômico-Decisórios & 4 & 2,7\% \\
\midrule
\textbf{Total} & \textbf{150} & \textbf{100\%} \\
\bottomrule
\end{tabular}
\end{table}

O Bloom total para 150 questões é: \bloom{Aplicar}=87 (58,0\%) e
\bloom{Analisar}=63 (42,0\%).

% ============================================================
\section{Análise de Correlações Estatísticas}
% ============================================================

Foram calculadas três matrizes de correlação utilizando o teste de
$\chi^2$ (qui-quadrado) e o V de Cramer como medida de associação.
O V de Cramer varia de 0 (associação nula) a 1 (associação perfeita).

\subsection{Área da Física $\times$ R-code Primário}

\begin{table}[H]
\centering
\caption{Tabela de contingência: Área $\times$ R-code primário}
\label{tab:area_rcode}
\begin{tabular}{lcccccccccccccc}
\toprule
\textbf{Área} & R008 & R010 & R011 & R065 & R067 & R069 & R070 & R071 & R073 & R109 & R117 & R183 & R201 & R202 & R229 \\
\midrule
Eletromagnetismo & 0 & 0 & 0 & 3 & 0 & 0 & 0 & 9 & 0 & 0 & 0 & 0 & 0 & 0 & 0 \\
Física Moderna & 0 & 0 & 0 & 2 & 0 & 0 & 1 & 1 & 0 & 0 & 0 & 8 & 0 & 0 & 0 \\
Mec. Clássica & 0 & 0 & 1 & 3 & 0 & 7 & 0 & 0 & 0 & 0 & 0 & 0 & 0 & 0 & 1 \\
Mec. Estatística & 0 & 0 & 0 & 0 & 0 & 1 & 0 & 0 & 1 & 0 & 4 & 0 & 0 & 0 & 0 \\
Mec. Quântica & 0 & 1 & 0 & 0 & 1 & 0 & 0 & 0 & 0 & 1 & 0 & 3 & 1 & 2 & 3 \\
Termodinâmica & 1 & 0 & 0 & 0 & 0 & 0 & 0 & 0 & 0 & 0 & 5 & 0 & 0 & 0 & 0 \\
\midrule
\textbf{Total} & 1 & 1 & 1 & 8 & 1 & 8 & 1 & 10 & 1 & 1 & 9 & 11 & 1 & 2 & 4 \\
\bottomrule
\end{tabular}
\end{table}

\begin{itemize}
  \item $\chi^2 = 218,8$ (gl = 70)
  \item \textbf{V de Cramer = 0,770} (associação \textbf{muito forte})
  \item $p < 10^{-6}$
\end{itemize}

\textbf{Interpretação:} A associação muito forte entre área e R-code
confirma que cada área da física mobiliza um repertório característico
de raciocínios. Os R-codes dominantes por área são:
\begin{itemize}
  \item Eletromagnetismo: \rcode{R071} (Dualidade-Transformacao) --- 9/12
  \item Física Moderna: \rcode{R183} (Mecanica-Quantica-Particulas) --- 8/12
  \item Mecânica Clássica: \rcode{R069} (Variacional-Principio) --- 7/12
  \item Mecânica Estatística: \rcode{R117} (Monte-Carlo-MCMC) --- 4/6
  \item Mecânica Quântica: \rcode{R183/R229/R202/R201} --- diversificada
  \item Termodinâmica: \rcode{R117} (Monte-Carlo-MCMC) --- 5/6
\end{itemize}

\subsection{Nível de Bloom $\times$ R-code Primário}

\begin{table}[H]
\centering
\caption{Tabela de contingência: Bloom $\times$ R-code primário}
\label{tab:bloom_rcode}
\begin{tabular}{lcccccccccccccc}
\toprule
\textbf{Bloom} & R008 & R010 & R011 & R065 & R067 & R069 & R070 & R071 & R073 & R109 & R117 & R183 & R201 & R202 & R229 \\
\midrule
Aplicar & 0 & 1 & 1 & 5 & 0 & 6 & 0 & 7 & 0 & 1 & 5 & 8 & 1 & 1 & 1 \\
Analisar & 1 & 0 & 0 & 3 & 1 & 2 & 1 & 3 & 1 & 0 & 4 & 3 & 0 & 1 & 3 \\
\midrule
\textbf{Total} & 1 & 1 & 1 & 8 & 1 & 8 & 1 & 10 & 1 & 1 & 9 & 11 & 1 & 2 & 4 \\
\bottomrule
\end{tabular}
\end{table}

\begin{itemize}
  \item $\chi^2 = 22,5$ (gl = 14)
  \item \textbf{V de Cramer = 0,602} (associação \textbf{forte})
  \item $p < 0,07$ (significância marginal)
\end{itemize}

\textbf{Interpretação:} A associação forte (V = 0,602) entre Bloom e
R-code sugere que alguns raciocínios tendem a ser mobilizados em níveis
cognitivos específicos. R-codes como \rcode{R008} (Deducao-Formal),
\rcode{R070} (Limite-Assintotico) e \rcode{R073} (Caos-Deterministico)
aparecem exclusivamente como \bloom{Analisar}, enquanto \rcode{R010}
(Decomposicao-Modular), \rcode{R011} (Inducao-Estrutural) e \rcode{R201}
(Superposicao-Quantica) aparecem exclusivamente como \bloom{Aplicar}.

\subsection{Área da Física $\times$ Nível de Bloom}

\begin{table}[H]
\centering
\caption{Tabela de contingência: Área $\times$ Bloom}
\label{tab:area_bloom}
\begin{tabular}{lccc}
\toprule
\textbf{Área} & \textbf{Aplicar} & \textbf{Analisar} & \textbf{Total} \\
\midrule
Eletromagnetismo & 10 & 2 & 12 \\
Física Moderna & 4 & 8 & 12 \\
Mecânica Clássica & 10 & 2 & 12 \\
Mecânica Estatística & 4 & 2 & 6 \\
Mecânica Quântica & 6 & 6 & 12 \\
Termodinâmica & 3 & 3 & 6 \\
\midrule
\textbf{Total} & \textbf{37} & \textbf{23} & \textbf{60} \\
\bottomrule
\end{tabular}
\end{table}

\begin{itemize}
  \item $\chi^2 = 5,2$ (gl = 5)
  \item \textbf{V de Cramer = 0,232} (associação \textbf{fraca})
  \item $p < 0,39$ (não significativo)
\end{itemize}

\textbf{Interpretação:} A associação fraca entre área e Bloom indica que
o nível cognitivo das questões não é determinado pela área da física, mas
sim pelo tipo de raciocínio mobilizado. A distribuição de Bloom é
relativamente homogênea entre as áreas, com exceção de Eletromagnetismo
e Mecânica Clássica (predominantemente \bloom{Aplicar}) e Física Moderna
(predominantemente \bloom{Analisar}).

\subsection{Matriz de Correlações}

A Tabela~\ref{tab:matriz_corr} resume as três correlações calculadas.

\begin{table}[H]
\centering
\caption{Matriz de correlações: $\chi^2$, V de Cramer e interpretação}
\label{tab:matriz_corr}
\begin{tabular}{lcccc}
\toprule
\textbf{Correlação} & $\chi^2$ & \textbf{gl} & \textbf{V de Cramer} & \textbf{Interpretação} \\
\midrule
Área $\times$ R-code & 218,8 & 70 & 0,770 & Muito forte \\
Bloom $\times$ R-code & 22,5 & 14 & 0,602 & Forte \\
Área $\times$ Bloom & 5,2 & 5 & 0,232 & Fraca \\
\bottomrule
\end{tabular}
\end{table}

A interpretação conjunta revela uma hierarquia de determinação:
\begin{enumerate}
  \item A \textbf{área da física} determina fortemente o \textbf{R-code}
    mobilizado (V = 0,770);
  \item O \textbf{R-code} determina moderadamente o \textbf{nível de Bloom}
    (V = 0,602);
  \item A \textbf{área da física} tem efeito fraco sobre o \textbf{Bloom}
    (V = 0,232).
\end{enumerate}

Isto sugere um modelo de mediação: o efeito da área sobre o nível
cognitivo é indireto, sendo mediado pelo tipo de raciocínio mobilizado.
Em termos pedagógicos, a escolha do R-code (que tipo de raciocínio
será exigido) é o principal determinante da profundidade cognitiva da
questão.

% ============================================================
\section{R-codes Dominantes por Área}
% ============================================================

\subsection{R071 --- O R-code Universal}

O \rcode{R071} (Dualidade-Transformacao) é o único R-code presente como
primário em \textbf{todos os 6 exames}, com 10 ocorrências totais. Sua
ubiquidade reflete a importância central das transformações matemáticas
(transformações de Lorentz, transformadas de Fourier, mudanças de
coordenadas) em todas as áreas da física.

\subsection{R069 e R065 --- O Núcleo Metodológico da Física}

\rcode{R069} (Variacional-Principio) e \rcode{R065}
(Simetria-Conservacao) aparecem em 5 dos 6 exames, com 8 ocorrências
cada. Eles formam uma dupla complementar: identificar simetrias $\to$
formular princípio variacional $\to$ extrair equações de movimento.

\subsection{R183 e R117 --- Especialização por Área}

\rcode{R183} (Mecanica-Quantica-Particulas) é dominante em Física
Moderna (8/12) e \rcode{R117} (Monte-Carlo-MCMC) domina Termodinâmica
(5/6) e Mecânica Estatística (4/6). Estes R-codes representam
raciocínios especializados, típicos de subdisciplinas específicas.

\subsection{R-codes com Ocorrência Única}

Seis R-codes aparecem como primário em apenas 1 das 60 questões:
\rcode{R008}, \rcode{R010}, \rcode{R011}, \rcode{R067}, \rcode{R070},
\rcode{R073}, \rcode{R201}. Eles representam raciocínios de fronteira,
mobilizados apenas em questões específicas que exigem abordagens
incomuns.

% ============================================================
\section{Comparação EUF 2014-2 com os Demais Exames}
% ============================================================

\subsection{Diferenças Qualitativas}

A EUF 2014-2 se distingue dos demais exames em três aspectos principais:

\begin{enumerate}
  \item \textbf{Maior diversidade categórica:} 4 categorias diferentes
    (I, II, XIII, XIV), contra 3 categorias nos demais exames;
  \item \textbf{Ausência de R-codes quânticos:} \rcode{R183}
    (Mecanica-Quantica-Particulas) está ausente como primário, sendo
    substituído por \rcode{R109} (Espectral-Operadores) e \rcode{R010}
    (Decomposicao-Modular) nas questões de Mecânica Quântica;
  \item \textbf{R-codes exclusivos:} \rcode{R008}, \rcode{R010},
    \rcode{R070} e \rcode{R109} como primários só aparecem na EUF 2014-2.
\end{enumerate}

\subsection{Similaridades}

Todos os exames compartilham:
\begin{itemize}
  \item Predomínio da Cat.~XIII (Físico-Matemáticos) com 7 a 9 ocorrências;
  \item Presença de \rcode{R071} (Dualidade-Transformacao) em todos;
  \item Distribuição Bloom aproximadamente 6A/4AN;
  \item Estrutura de 10 questões cobrindo as mesmas 6 áreas.
\end{itemize}

% ============================================================
\section{Conclusões e Implicações}
% ============================================================

A análise completa dos 6 exames EUF permite concluir que:

\begin{enumerate}
  \item \textbf{A Taxonomia 350 é um instrumento eficaz} para classificar
    o raciocínio científico em exames de física, discriminando 20 tipos
    de raciocínio em 5 categorias;

  \item \textbf{O V de Cramer Área $\times$ R-code = 0,770} valida
    quantitativamente a hipótese de que cada área da física tem um
    perfil característico de raciocínio;

  \item \textbf{O nível de Bloom é mediado pelo R-code}, não diretamente
    pela área, sugerindo que o projeto pedagógico de questões deve
    focar na escolha do tipo de raciocínio;

  \item \textbf{A EUF 2014-2 é o exame mais diverso} do dataset em
    termos de categorias taxonômicas, sendo o único a mobilizar
    raciocínios lógico-dedutivos (Cat.~I) e analítico-estruturais
    (Cat.~II) como primários;

  \item \textbf{O repertório de R-codes se expandiu} nos exames mais
    recentes (2015-2, 2016-1), que incorporam R-codes de mecânica
    quântica avançada (R201, R202) ausentes nos exames anteriores.
\end{enumerate}

\subsection{Trabalhos Futuros}

A continuidade desta pesquisa pode incluir:
\begin{itemize}
  \item Análise de correlação entre desempenho dos candidatos e tipos
    de raciocínio;
  \item Modelagem preditiva do nível de Bloom a partir do R-code e da
    área;
  \item Validação inter-avaliadores da classificação taxonômica;
  \item Extensão para outras disciplinas científicas (Química, Biologia);
  \item Análise de séries temporais da evolução dos perfis taxonômicos
    ao longo dos anos em ambos os exames.
\end{itemize}

\end{document}
